\documentclass[11pt,a4paper]{article}
\usepackage[utf8]{inputenc}
\usepackage{amsmath}
\usepackage{amsfonts}
\usepackage{graphicx}
\usepackage[table,xcdraw]{xcolor}

\DeclareUnicodeCharacter{00A0}{ } %este rollo es para evitar un error por la aparición de un caracter invisible e hijo de puta


\usepackage{caption}
\usepackage{subcaption}

\renewcommand{\familydefault}{\sfdefault} % cambiamos la fuente a una sans

\usepackage{float} % para que floten las imagenes o algo asi...
\usepackage{wallpaper} %paquete para usar una imagen como encabezado!
\usepackage{hyperref} %para usar hypervinculos 
\usepackage[export]{adjustbox} %para usar marcos en imagenes
\usepackage{eurosym} % para el euro
\usepackage{transparent} %para las marcas de agua
\usepackage{eso-pic}  %para las marcas de agua

\definecolor{azul_marcos}{RGB}{0,128,159} %defino el color azul de los marcos
\usepackage{sectsty} %esto es para cambiar el color de las fuentes creo
\renewcommand{\familydefault}{\sfdefault} % cambiamos la fuente a una sans
\sectionfont{\color{azul_marcos}}  % sets colour of sections
\subsectionfont{\color{azul_marcos}}  % sets colour of sections


\usepackage{pdfpages} %para insertar pdfs
\usepackage{amssymb}
\usepackage{pstcol} % para color
\usepackage{pst-node} % para diagramas
\usepackage{pst-plot} % para representacion de dat
\usepackage[spanish]{babel}
\addto\captionsspanish{\renewcommand\chaptername{Bloque}}
%\usepackage[total={18cm,21cm},top=2cm, left=2cm]{geometry}
\usepackage{anysize}
\pagestyle{plain}
%\markboth{left head}{right head}
%\markright{Guía de impresión ABS Tech}
\marginsize{3cm}{2cm}{2.5cm}{1cm}
\title{ABS Tech printing guide}
\date{}

%configuracion de la marca de agua
\AddToShipoutPicture{
    \put(0,0){
        \parbox[b][\paperheight]{\paperwidth}{%
            \vfill
            \centering
            {\transparent{0.1}\includegraphics[scale=1.25]{FOTOS/logofff}}%
            \vfill
        }
    }
}

\begin{document}
\ULCornerWallPaper{1}{FOTOS/header}
\LLCornerWallPaper{1}{FOTOS/footer}
%\maketitle
%\tableofcontents
\includepdf{PDF/PORTADA.pdf}
\section{What is the ABS Tech?}ABS Tech is a filament for 3D printing FFF/FDM of acrylonitrile butadiene styrene (A.B.S.) specially treated with additives to reduce the warping effect which is characteristic of this material and also to make it easier to use.
\section{Why should you use ABS Tech?}
ABS Tech expands the possibilities of the traditional ABS filaments while maintaining at the same time the chemical and mechanical properties of this great thermoplastic.
\begin{itemize}
\item Its anti-warping treatment enables the printing in larger volumes than the standard and a better use in the household printers.
\item The odors produced by our ABS Tech are less annoying than those of other filaments of ABS Tech.
\item ABS Tech is a very tough material and with good chemical resistance to abrasion.
\item It dissolves in acetone. This allows to join pieces of ABS tightly using acetone as glue.
\item And finally the printed parts in ABS Tech can be post-processed with ease as this material can be drilled, sanded, painted, etc...
\end{itemize}
\section{Technical data and the printing parameters}
\begin{table}[H]
\centering
\caption*{Technical Data}
\begin{tabular}{|
>{\columncolor[HTML]{FFFFFF}}l |
>{\columncolor[HTML]{FFFFFF}}c |}
\hline
\multicolumn{1}{|c|}{\cellcolor[HTML]{FFFFFF}\textbf{Material}}   & Acrylonitrile Butadiene Styrene   \\ \hline
\textbf{Available colours}              & 22                 \\ \hline
\textbf{Available Formats}             & 1kg, 250gr         \\ \hline
\textbf{Temperature of thermical deflection} & 88ºC               \\ \hline
\textbf{Fusion temperature}            & 200ºC              \\ \hline
\textbf{Decomposition temperature}    & \textgreater 260ºC \\ \hline
\textbf{Density}                         & 1.05 gr / cm3      \\ \hline
\textbf{Impact resistance}                         & 17 kg-cm / cm3      \\ \hline
\textbf{Maximum Stretching}              & 25\%              \\ \hline
\end{tabular}
\end{table}
\begin{table}[H]
\centering
\begin{tabular}{|
>{\columncolor[HTML]{FFFFFF}}l |
>{\columncolor[HTML]{FFFFFF}}c |}
\hline
\multicolumn{1}{|c|}{\cellcolor[HTML]{FFFFFF}\textbf{Recommended printing temperature}} & 240º-245º              \\ \hline
\textbf{Recommended printing speed}                         & 50-90mm/s              \\ \hline
\textbf{Heated bed temperature}                                  &  80º - 100º        \\ \hline
\textbf{Layer fan}                                      & Desactivado o a baja velocidad                 \\ \hline
\textbf{First layer speed}                                                 & 20 mm/s                      \\ \hline
\textbf{First layer heigth}                                          & \textgreater 0.2 mm                      \\ \hline
\end{tabular}
\end{table}

You can download our complete printing profiles of major programs of lamination (Cure, Slic3r and SImplify3D) from our website:
\\\\
\centerline{ {\huge \url{www.fffworld.com/documentation} } }
\\\\
The optimal parameters depend on the 3D printer that you use, however, they are good parameters for keeping them as a starting point. With a few prints you will be able to find the limits and the perfect setting for your machine.
\section{Technical Data}
	\subsection{The warping and the cracking}
		\subsubsection{What are they?}The ABS is the ideal material for 3D printing for their availability and the characteristics: it is extruded at a temperature slightly higher than the PLA, it is less demanding on the type of utilized hot-end, and it is less prone to jamming since it does not crystallize upon degradation and it possesses superior mechanical properties.
\\\\
Its biggest disadvantage is the known effect of warping / and cracking, however with the adequate tips these problems can be overcome.
\\\\
The warping effect is the name given in the world of 3D printing for the problems caused by plastic contraction upon cooling which sometimes causes the pieces to deform or to break.
\\\\
Distinctions are made between warping and cracking in accordance with the problem affecting the first layer or the intermediate layers of the part.
%EJEMPLO_WARPING
%Leyenda: Ejemplo de problema de warping
%EJEMPLO_CRACKING
%Leyenda: Ejemplo de problema de cracking
\begin{figure}[H]
    \centering
    \begin{subfigure}[b]{0.4\textwidth}
        \includegraphics[width=\textwidth,cfbox=azul_marcos 4pt 0pt]{FOTOS/EJEMPLO_WARPING}
	\caption*{Warping example}
    \end{subfigure}
    \qquad %add desired spacing between images, e. g. ~, \quad, \qquad, \hfill etc. 
      %(or a blank line to force the subfigure onto a new line)
    \begin{subfigure}[b]{0.4\textwidth}
        \includegraphics[width=\textwidth,cfbox=azul_marcos 4pt 0pt]{FOTOS/EJEMPLO_CRACKING}
	\caption*{Cracking example}
    \end{subfigure}   
\end{figure}
		\subsubsection{What are their causes?}Printing in 3D means depositing threads of filament that are glued together and they construct the desired pieces.
\\\\
Upon cooling these threads shrink reducing its length and causing cumulative stresses on the pieces with undesired effects.
\\\\
The different phases of the problem can be seen in the image:
%CAUSA WARPING
\begin{figure}[H]
\centering
\includegraphics[width=0.6\textwidth,cfbox=azul_marcos 4pt 0pt]{FOTOS/CAUSA_WARPING_1}
\end{figure}

\begin{description}
\item[1] The first layer is deposited hot: Being still warm its size is greater than that you will get at room temperature.
\item[2] The first layer cools: Upon cooling it shrinks its size and the first tensions appear which tend to detach the part of the printing surface.
\item[3] The second layer is deposited hot: when hot, the volume of the second layer is greater than the first.
\item[4] The second layer cools: The tension generated by the contraction of the second layer adds to that caused by the first layer and becomes detached while curving.
\end{description}
			\paragraph{The influence of the temperature in warping}\mbox{}\\\\
The temperature difference is thus responsible for the problem, particularly the difference between room temperature and the glass transition temperature of the material.
\\\\
\url{https://en.wikipedia.org/wiki/Glass_transition}
\\\\
The glass transition temperature of ABS is about 100º and the room temperature usually around 30º, is that jump of 70º which causes the problem.
\\\\
As a curiosity it is worth explaining that the PLA expands more than the ABS with temperature, nevertheless its glass transition stands at 60 degrees and since the difference with room temperature is much smaller, it is not affected by the warping.
\\\\
Therefore by controlling the room temperature we can fully control the warping, although this is not always easy.
		\subsubsection{Physical Solutions}
			\paragraph{Improving the adhesion of the printing surface}\mbox{}\\\\
The warping effect can be attenuated by improving the adhesion of the printing surface, however this is not the best solution because it does not attack the root of the problem and would not do anything to solve the cracking.
\\\\
%CAUSA WARPING 2
\begin{figure}[H]
\centering
\includegraphics[width=0.6\textwidth,cfbox=azul_marcos 4pt 0pt]{FOTOS/CAUSA_WARPING_2}
\end{figure}

\begin{description}
\item[1] The first layer is deposited hot: By being still hot its size is greater than that you will get at room temperature.
\item[2] The first layer cools: Upon cooling the first layer tends to shrink, however the adhesion of the surface counteracts this force.
\end{description}
By applying an adhesive product the tensions produced by cooling of the ABS will not be removed, but they can be counteracted and it is an effective solution for small pieces.
\\\\
It is very important that the surface is free of dust, grease or foreign elements that negatively affect the adhesion.
\\\\
It is also necessary that the printing surface is perfectly leveled so that the first layer has a uniform altitude.
\\\\
It is also necessary that the printing surface is perfectly leveled so that the first layer has a uniform altitude.

				\subparagraph{Kapton tape}\mbox{}\\\\
This is an adhesive tape with silicone adhesive polyamide prepared to withstand high temperatures. The ABS is glued very well to this type of tape and has been widely used in 3D printing with this purpose. In order to use it properly you have to stick it carefully to the printing surface.
%KAPTON1 KAPTON2
\begin{figure}[H]
    \centering
    \begin{subfigure}[b]{0.4\textwidth}
        \includegraphics[width=\textwidth,cfbox=azul_marcos 4pt 0pt]{FOTOS/KAPTON1}
	\caption*{Kapton tape}
    \end{subfigure}
    \qquad %add desired spacing between images, e. g. ~, \quad, \qquad, \hfill etc. 
      %(or a blank line to force the subfigure onto a new line)
    \begin{subfigure}[b]{0.4\textwidth}
        \includegraphics[width=\textwidth,cfbox=azul_marcos 4pt 0pt]{FOTOS/KAPTON2}
	\caption*{Cristal with Kapton}
    \end{subfigure}   
\end{figure}
				\subparagraph{The Lacquer}\mbox{}\\\\
The lacquer has proven to be an effective recourse in increasing the adhesion of the printing surface. It is recommended to choose a lacquer containing the least possible additives. The way to use it is to spray abundantly on the surface before printing.
%LACA
\begin{figure}[H]
\centering
\includegraphics[width=0.3\textwidth,cfbox=azul_marcos 4pt 0pt]{FOTOS/LACA}
\caption*{Nelly power lacquer}
\end{figure}
				\subparagraph{The dissolution of ABS-Acetone}\footnote{WARNING: This method can be so effective that the resulting part cannot be detached from the base without breaking one of the two.} \mbox{}\\\\
By taking advantage of the solubility of ABS in acetone you could prepare a solution to tightly glue the pieces of the printing surface. This is done by immersing the pieces of ABS in acetone and letting it dissolve the plastic. Once dissolved, the resulting liquid can be applied with a brush on the printing surface to greatly increase the adhesion.
				\subparagraph{Other products}\mbox{}\\\\
In addition to the methods mentioned there are users who use others ways to achieve the same goals. On the other side the advancement of 3D printing has encouraged the appearance of specific products in order to solve the adhesion problems? We believe that the above methods are the most affordable and effective interventions, however we want to mention some others because they can be of a great help:
\begin{description}
\item[Printing surfaces] 3M BlueTape, Buildtak
\item[Productos adhesivos] Fixwarp3d, 3DLAC, Dimafix, UHU Glue Stick(PVA)
\end{description}
			\paragraph{The use of heated bed}\mbox{}\\\\
Using a heated printing surface (heated bed) is required to print ABS in a satisfactory manner and there are many printers that incorporate one of this.
\begin{figure}[H]
\centering
\includegraphics[width=0.6\textwidth,cfbox=azul_marcos 4pt 0pt]{FOTOS/HEATEDBED}
\caption*{RepRap printer heated bed}
\end{figure}
The hot bed will ensures that the first layers are maintained at a temperature sufficient in order to avoid the material shrinkage and the problems of warping.
\\\\
The hot bed should be heated at least 80º although it is advisable that it should be at least 100º.
\\\\
Must be taken into account that the hot bed does not get to keep the temperature of the upper layers in large pieces therefor it may still occur the cracking and especially in large volumes.
			\paragraph{Avoid the air currents}\mbox{}\\\\
In order to prevent that the part can be cooled down abruptly it is very convenient to print in the areas isolated from the air currents. This enters in conflict with the recommendation of printing in the ventilated areas in order to prevent the accumulation of harmful gases.
\\\\
Therefore the best option is to enclose the printer in an enclosure capable of maintaining the temperature and to avoid the air currents.
% CERRAMIENTO1  CERRAMIENTO2  CERRAMIENTO3 
\begin{figure}[H]
    \centering
    \begin{subfigure}[b]{0.3\textwidth}
        \includegraphics[width=\textwidth,cfbox=azul_marcos 4pt 0pt]{FOTOS/CERRAMIENTO1}
    \end{subfigure}
    \quad %add desired spacing between images, e. g. ~, \quad, \qquad, \hfill etc. 
      %(or a blank line to force the subfigure onto a new line)
    \begin{subfigure}[b]{0.3\textwidth}
        \includegraphics[width=\textwidth,cfbox=azul_marcos 4pt 0pt]{FOTOS/CERRAMIENTO2}
    \end{subfigure}
    \quad %add desired spacing between images, e. g. ~, \quad, \qquad, \hfill etc. 
    %(or a blank line to force the subfigure onto a new line)
    \begin{subfigure}[b]{0.3\textwidth}
        \includegraphics[width=\textwidth,cfbox=azul_marcos 4pt 0pt]{FOTOS/CERRAMIENTO3}
    \end{subfigure}
    \caption*{Homemade enclosure examples}
\end{figure}
Meanwhile there are many printers that are designed in the shape of the box and do not need this additional insulation.
% IMPRESORACERRADA1	IMPRESORACERRADA2	IMPRESORACERRADA3
\begin{figure}[H]
    \centering
    \begin{subfigure}[b]{0.3\textwidth}
        \includegraphics[width=\textwidth,cfbox=azul_marcos 4pt 0pt]{FOTOS/IMPRESORACERRADA1}
    \end{subfigure}
    \quad %add desired spacing between images, e. g. ~, \quad, \qquad, \hfill etc. 
      %(or a blank line to force the subfigure onto a new line)
    \begin{subfigure}[b]{0.3\textwidth}
        \includegraphics[width=\textwidth,cfbox=azul_marcos 4pt 0pt]{FOTOS/IMPRESORACERRADA2}
    \end{subfigure}
    \quad %add desired spacing between images, e. g. ~, \quad, \qquad, \hfill etc. 
    %(or a blank line to force the subfigure onto a new line)
    \begin{subfigure}[b]{0.3\textwidth}
        \includegraphics[width=\textwidth,cfbox=azul_marcos 4pt 0pt]{FOTOS/IMPRESORACERRADA3}
    \end{subfigure}
    \caption*{Enclosed 3d printers}
\end{figure}
			\paragraph{Heating the enclosure}\mbox{}\\\\
The industrial and professional printers print ABS in a heated enclosure that keeps the part at 80º during all the printing. This requires extra engineering to cool down the hot-end and the rest of the components, but it completely eliminates the problems of warping and cracking.
\\\\
% STRATASYS
\begin{figure}[H]
\centering
\includegraphics[width=0.6\textwidth,cfbox=azul_marcos 4pt 0pt]{FOTOS/STRATASYS}
\caption*{Impresora industrial con recinto calefactado}
\end{figure}
Although this kind of printers are not available to the domestic user we will mention them for a better understanding of the problem.
		\subsubsection{Software and design solutions}Besides all the above mentioned techniques they also do exist other design and lamination techniques that allow to reduce or eliminate the problems of warping.
			\paragraph{Brim and raft}\mbox{}\\\\
The brim and raft are options that offer programs of lamination to increase the adhesion of the pieces to the base.
%BRIM			RAFT 	
% Leyenda: Brim	Leyenda: Raft
\begin{figure}[H]
    \centering
    \begin{subfigure}[b]{0.4\textwidth}
        \includegraphics[width=\textwidth,cfbox=azul_marcos 4pt 0pt]{FOTOS/BRIM}
    \end{subfigure}
    \qquad %add desired spacing between images, e. g. ~, \quad, \qquad, \hfill etc. 
      %(or a blank line to force the subfigure onto a new line)
    \begin{subfigure}[b]{0.4\textwidth}
        \includegraphics[width=\textwidth,cfbox=azul_marcos 4pt 0pt]{FOTOS/RAFT}
    \end{subfigure}   
\end{figure}
The Brim option creates an extra perimeter around the piece to increase the surface and to improve the adhesion.
\\\\
Meanwhile the Raft option creates a kind of bed of several layers of thickness and then print that piece on this bed.
\\\\
Generally it is preferable to use Brim because it is faster.
\\\\
For more information about how to activate and use these options in your program you can see the following:\\\\
\url{http://manual.slic3r.org/expert-mode/skirt}\\\\
\url{https://www.simplify3d.com/support/tutorials/rafts-skirts-and-brims/}\\\\
\url{https://ultimaker.com/en/resources/16525-platform-adhesion}
			\paragraph{The Impact of infill in warping}\mbox{}\\\\
The degree of warping is directly proportional to the amount of plastic shrink, thus the lower percentages of the infill will produce less warping than those pieces with more filling.
\\\\
The heigth of the first layers also affects and may be better to get it reduced.
			\paragraph{The modification of parts for reducing warping}
				\subparagraph{Add the supports around the corners}\mbox{}\\\\
A method of controlling the warping is to redesign the piece by strengthening the points where it has been detached from the bed.
\\\\
If after a failed printing we check to see that one or more corners continue rising without remedy and it may be necessary to add a support in this area.
\\\\
The support must be larger or smaller depending on the severity of the problem.
%	MOUSEEAR1			MOUSEEAR2
\begin{figure}[H]
    \centering
    \begin{subfigure}[b]{0.4\textwidth}
        \includegraphics[width=\textwidth,cfbox=azul_marcos 4pt 0pt]{FOTOS/MOUSEEAR1}
    \end{subfigure}
    \qquad %add desired spacing between images, e. g. ~, \quad, \qquad, \hfill etc. 
      %(or a blank line to force the subfigure onto a new line)
    \begin{subfigure}[b]{0.4\textwidth}
        \includegraphics[width=\textwidth,cfbox=azul_marcos 4pt 0pt]{FOTOS/MOUSEEAR2}
    \end{subfigure}   
\end{figure}	
				\subparagraph{Avoid the rounded corners}\mbox{}\\\\
It has been proven that the warping affects more round corners and the parts with circular base or convex perimeters.
\\\\
If you are having this problem your part can benefit from a redesign that eliminates such geometry problems.
			\paragraph{Lowering the speed and the height of the first layer}\mbox{}\\\\
It is very important that the first layer stays as best as possible adhered to the printing surface.
\\\\
To achieve this it is advisable to significantly reduce the speed of the first layer by favoring that the material sticks in a more firm and uniform way at the base.
\\\\
A speed of 20 mm /s in the first layer should be sufficient to achieve that objective. It is also advisable to reduce the height of the first layer so that does not exceed more than 0.2 mm.
			\paragraph{Using of layer fan}\mbox{}\\\\
The ventilator of the layer is responsible for cooling the plastic after its extrusion and this is precisely what to avoid when printing ABS.
\\\\
Therefore as a general rule it is recommended to disable it while printing ABS.
\\\\
However it can be used selectively in some moments of printing if your laminated software allows it.
\\\\
However it can be used selectively in some moments of printing if your laminated software allows it.
	\subsection{The toxic vapors}ABS is known to emit harmful vapors when printed which it can be potentially harmful to humans.
\\\\
So it is recommended that the printer is in a ventilated place or at least do not stay with it for a long time during the printing process.
\\\\
If you use intensively ABS it may be interesting to install a filtering and extracting system in your printer to completely avoid being exposed to these vapors.
% VENTILACIONVAPORES
\begin{figure}[H]
\centering
\includegraphics[width=0.5\textwidth,cfbox=azul_marcos 4pt 0pt]{FOTOS/VENTILACIONVAPORES}
\caption*{Vapors ventilation}
\end{figure}
\section{Printing Tips: The post processing of ABS}
	\subsection{Sanding, cutting and drilling}One of the advantages of ABS is its ease of the post processing. The pieces of ABC can be sanded with a conventional sandpaper to remove the marks of the printing process or to facilitate its assembly.
\\\\
It can also be drilled without any problems and it can be shaped by cutting with a cutter or something similar.
	\subsection{The paint}The ABS can be painted by using acrylic paint. It is convenient to sand the surface of the workpiece for a better grip of the paint.
	\subsection{Acetone as glue}Being soluble in acetone it may be used as a glue to strongly join pieces of ABS. By applying a small amount on the contact surfaces they will be merged with each other by providing a strong and durable bond. 
	\subsection{Acetone for polishing pieces}The acetone can also be used as a surface treatment to give the pieces of ABS a perfect polished and shiny look.
% VAPORACETONA1			VAPORACETONA2
\begin{figure}[H]
    \centering
    \begin{subfigure}[b]{0.4\textwidth}
        \includegraphics[width=\textwidth,cfbox=azul_marcos 4pt 0pt]{FOTOS/VAPORACETONA1}
    \end{subfigure}
    \qquad %add desired spacing between images, e. g. ~, \quad, \qquad, \hfill etc. 
      %(or a blank line to force the subfigure onto a new line)
    \begin{subfigure}[b]{0.4\textwidth}
        \includegraphics[width=\textwidth,cfbox=azul_marcos 4pt 0pt]{FOTOS/VAPORACETONA2}
    \end{subfigure}   
\end{figure}
It can be applied by hand using a brush, but it is also possible to create relatively easily a home steam chamber for the acetone. In the following links are described two different methods of bathing with acetone vapor the pieces of ABS:
\\\\
\url{http://www.instructables.com/id/Safe-way-to-do-Acetone-bath/}\\\\
\url{http://sinkhacks.com/building-acetone-vapor-bath-smoothing-3d-printed-parts/}
\section{Would you like to support our project?}
All the members of FFF World love 3D printing and the maker community. We feel lucky to be able to work on projects where we can deliver our honest passion. In the future, we would like to be able to develop more materials, more colors and more formats. Ultimately, we would like to be able to make our company grow.
\\\\
Therefore, one of the best actions to help us, if you want to do it and you’re satisfied with the filament, is to give us a 5 stars rating on amazon.
\begin{figure}[H]
\centering
\includegraphics[width=0.5\textwidth,cfbox=azul_marcos 1pt 0pt]{FOTOS/AMAZON_FIVE_STARS}
\caption*{Thanks a lot!}
\end{figure}
\subsection{Other filaments with awesome properties available today in Amazon}
\begin{description}
\item[FlexiSMART Tech:] Designed to resist the abrasion and wear of technical printings.
\item[ABS Tech:] Minimized warping effect. High performance on technical applications.
\item[PETG Tech:] Maximum mechanical resistance. Resistant to contact with water and UV rays. Apt for alimentary use.
\item[FilaMETAL:] PLA with non-abrasive metallic charge that gives a spectacular mechanical finish to your printings.
\item[PC Tech:] Polycarbonates with a high resistance to temperature and excellent mechanic properties.
\item[Nylon Tech:] printable at low temperature. Resistance to bangs with a certain degree of flexibility.
\item[PVA Tech:] Water soluble filament indicated for use as support material. Excellent compatibility with PLA.
\item[HIPS Tech:] Limonene soluble filament indicated for use as support material. Good mechanical resistance and excellent compatibility with ABS.
\end{description}
%\section{Bibliografía}
%Esta guía no habría sido posible sin el conocimiento libre generado por la comunidad RepRap. Para la elaboración de esta guía se han %utilizado imágenes y contenido extraidos de los siguientes sitios web.
%\\\\
%\url{http://www.gyrobot.co.uk/blog/how-to-3d-print-with-flexible-filaments}\\
%\url{http://www.thingiverse.com/thing:1496895}\\
%\url{http://www.thingiverse.com/thing:247024}\\
%\url{http://www.thingiverse.com/thing:16319}\\
%\url{http://www.thingiverse.com/thing:779011}\\
%\url{http://www.thingiverse.com/thing:1102900}\\
%\url{http://www.thingiverse.com/thing:147705}\\
%\url{http://www.thingiverse.com/thing:222667}\\
%\url{http://www.thingiverse.com/thing:512338}\\
%\url{https://all3dp.com/common-3d-printing-problems-and-their-solutions/}\\
%\url{https://www.simplify3d.com/support/}\\
%\url{http://www.thingiverse.com/thing:508896}\\
%\url{http://www.thingiverse.com/thing:1187344}
\includepdf{PDF/EN_CONTRAPORTADA.pdf}
\end{document}
